\section*{Hoofdstuk \ref{ch:b2:eruptive-styles} --- Vulkaanuitbarstingen}

\begin{solution}{exo:b2:eruptive-styles:1}
(a) $\qty{1.3}{km^3} = \qty{1.3e9}{m^3}$, tussen $10^9$ en $10^{10}$: VEI~5.
(b) $\qty{2.9e8}{m^3}$: VEI~4. (c) $\qty{2e6}{m^3}$, tussen $10^6$ en $10^7$:
VEI~2. (d) VEI~6 begint bij \qty{e10}{m^3}, dus de bulktefra moet minstens
\qty{10}{km^3} zijn; ``meer dan \qty{5}{km^3}'' is een ondergrens, geen tegenspraak.
\end{solution}

\begin{solution}{exo:b2:eruptive-styles:2}
(a) As (fijner dan \qty{2}{mm}). (b) Een lapillus (\qtyrange{2}{64}{mm}).
(c) Een blok: groter dan \qty{64}{mm}, hoekig, vast bij het uitwerpen. (d) Een
\omterm{def:b2:eruptive-styles:tephra}{vulkanische bom}, een spoelbom: groter dan \qty{64}{mm} en gevormd terwijl hij gesmolten was,
tijdens de vlucht.
\end{solution}

\begin{solution}{exo:b2:eruptive-styles:3}
(a) $0.304\times(10^6)^{0.241} = 0.304\times27.9 = \qty{8.5}{km}$;
$0.304\times(10^8)^{0.241} = 0.304\times84.7 = \qty{25.8}{km}$.
(b) $Q = (35/0.304)^{1/0.241} = 115.1^{4.149} = \qty{3.6e8}{kg/s}$, en $I =
8.55 + 3 = 11.6$. (c) $Q \propto H^{4.15}$, dus een hoogte die \qty{20}{\percent} te hoog is,
overschat $Q$ met een factor $1.2^{4.15} = 2.1$.
\end{solution}

\begin{solution}{exo:b2:eruptive-styles:4}
$V_{\mathrm{DRE}} = 0.29\times1000/2500 = \qty{0.116}{km^3}$; $m =
\num{0.29e9}\times1000 = \qty{2.9e11}{kg}$; $M = \log_{10}(\num{2.9e11}) - 7 = 11.46
- 7 = 4.46$.
\end{solution}

\begin{solution}{exo:b2:eruptive-styles:5}
$k = \ln(40/10)/(15 - 5) = 1.386/10 = \qty{0.139}{km^{-1}}$; $T_0 = 40\,e^{0.139\times5}
= 40\times2.00 = \qty{80}{cm}$; $\ln2/k = \qty{5.0}{km}$. Met $T_0 =
\qty{8.0e-4}{km}$: $V = 2\times\num{8.0e-4}/0.139^2 = \qty{0.083}{km^3} =
\qty{8.3e7}{m^3}$, VEI~3.
\end{solution}

\begin{solution}{exo:b2:eruptive-styles:6}
(a) $\int_0^x T_0e^{-ky}2y\,dy = 2T_0\big[-\tfrac{y}{k}e^{-ky} - \tfrac{1}{k^2}e^{-ky}\big]_0^x
= \tfrac{2T_0}{k^2}[1 - (1 + kx)e^{-kx}]$. (b) Iteratie van $u = \ln(2(1 + u))$ vanaf
$u = 1$: $1.386$, $1.564$, $1.635$, $1.662$, $1.672$, \dots, convergerend naar
$u = 1.678$, waarvoor $(1 + u)e^{-u} = 2.678\times0.1867 = 0.500$. (c) $x_{1/2} =
1.678/0.139 = \qty{12.1}{km}$, binnen de buitenste isopach ($\sqrt A =
\qty{15}{km}$); toch is het volume voorbij die isopach $(1 + 2.09)e^{-2.09} =
\qty{38}{\percent}$ van het totaal. Bijna twee vijfde van de afzetting ligt waar
niets werd gemeten: het volume berust op de extrapolatie.
\end{solution}

\begin{solution}{exo:b2:eruptive-styles:7}
(a) $p_s = (5.0/0.444)^2 = \qty{127}{MPa}$, op $127/0.025 = \qty{5.1}{km}$.
(b) $x = (5.0 - 0.444\sqrt{10})/100 = (5.0 - 1.404)/100 = 0.0360$;
$V_g = 0.0360\times8.314\times1123/(0.018015\times10^7) = \qty{1.86e-3}{m^3}$;
$V_m = 0.964/2300 = \qty{4.19e-4}{m^3}$; $\phi = 1.86/(1.86 + 0.419) = 0.82$.
(c) Een basaltsmelt is duizenden tot miljarden keren minder viskeus dan een
rhyolitische (\cref{ch:b2:magma-emplacement}): zijn bellen stijgen op,
versmelten en ontsnappen, en het gas vertrekt door een doorlatend schuim of bij de
krater, dus het \omterm{def:b2:magma-emplacement:magma}{magma} ontgast in een open systeem in plaats van zijn gas vast te houden; zijn
fragmentatie, als ze optreedt, blijft beperkt tot een fontein dicht bij de krater. Een tragere
opstijging laat hier meer tijd voor.
\end{solution}

\begin{solution}{exo:b2:eruptive-styles:8}
(a) $u_z(0) = 0.75\times10^7/(\pi\times(5000)^2) = \qty{0.095}{m} = \qty{95}{mm}$.
(b) $z_0 = 3.8/0.766 = \qty{5.0}{km}$. (c) $(1 + 1)^{-3/2} = 0.35$ van het
maximum.
\end{solution}

\begin{solution}{exo:b2:eruptive-styles:9}
De grootste dracht wordt bereikt bij \ang{45}: $R = v^2/g = 100^2/9.81 = \qty{1.0}{km}$, en
$200^2/9.81 = \qty{4.1}{km}$. Luchtweerstand verkort de vlucht, dus dit zijn
bovengrenzen, dicht bij de werkelijkheid voor grote blokken en veel te lang voor kleine.
Omdat $R = v^2/g$, groeit de gevarenzone met het kwadraat van de lanceersnelheid:
een verdubbeling van de snelheid van de explosie verviervoudigt haar straal.
\end{solution}

\begin{solution}{exo:b2:eruptive-styles:10}
(a) Met $T = T_0e^{-k_1x}$ voor $x < X$ en $T = T_Xe^{-k_2(x - X)}$ daarbuiten, $T_X =
T_0e^{-k_1X}$ (continuïteit), gebruik $\int_0^X ye^{-ky}dy = \frac{1}{k^2} - e^{-kX}\big(\frac{X}{k} +
\frac{1}{k^2}\big)$ en $\int_X^\infty ye^{-ky}dy = e^{-kX}\big(\frac{X}{k} +
\frac{1}{k^2}\big)$:
het proximale deel geeft $2T_0\big[\frac{1}{k_1^2} - e^{-k_1X}\frac{k_1X + 1}{k_1^2}\big]$ en het
distale deel $2T_Xe^{k_2X}\int_X^\infty ye^{-k_2y}dy = 2T_X\frac{k_2X + 1}{k_2^2}$; met $T_X =
T_0e^{-k_1X}$ is hun som de formule van het methodekader.
(b) $T_0 = \qty{0.01285}{km}$: $2T_0/k_1^2 = 0.0257/0.1332 = \qty{0.193}{km^3}$;
$e^{-0.365\times8.56} = 0.0440$; $(0.138\times8.56 + 1)/0.138^2 = 114.5$ en
$(0.365\times8.56 + 1)/0.365^2 = 30.96$; de correctie is $0.0257\times0.0440\times83.6
= \qty{0.094}{km^3}$ en $V = \qty{0.287}{km^3}$, de gepubliceerde \qty{0.29}{km^3}.
Voorbij $X$ ligt $2T_X(X/k_2 + 1/k_2^2) = 2\times\num{5.65e-4}\times114.5 =
\qty{0.129}{km^3}$, dat is \qty{45}{\percent}.
\end{solution}

\begin{solution}{exo:b2:eruptive-styles:11}
(a) $100^{1/4} = 3.16$; de empirische wetten geven $100^{0.241} = 3.03$ en
$100^{0.226} = 2.83$. (b) $H \propto N^{-3/4}$: een sterker gelaagde
atmosfeer stopt de kolom lager; voor een $N$ die \qty{20}{\percent} groter is,
$1.2^{-0.75} = 0.87$, een hoogte die \qty{13}{\percent} lager is. (c) Zwakke pluimen worden
door de wind omgebogen, wat ze verlaagt, terwijl vochtige lucht er latente warmte in
vrijmaakt, wat ze verhoogt; deze effecten wegen verschillend door bij kleine en grote
erupties en zitten niet in de eenvoudige wet. En de debieten zijn slechts tot op een factor
van enkele nauwkeurig bekend, een fout in de variabele op de horizontale as die
een gefitte helling naar nul vertekent.
\end{solution}

\begin{solution}{exo:b2:eruptive-styles:12}
(a) $132/400\times100 = 33$ per eeuw vóór 1900; $95/125\times100 = 76$ per
eeuw daarna. (b) $33/76 = \qty{43}{\percent}$. (c) $21/56 = \qty{38}{\percent}$:
dezelfde onderregistratie, met een factor twee tot drie. (d) Grotere erupties worden
van verder gezien en laten dikkere, uitgestrektere afzettingen na die de erosie
langer overleven; een eruptie van VEI~4 op een afgelegen eiland in 1700 kon onopgemerkt blijven, een
eruptie van VEI~6 nauwelijks.
\end{solution}

\begin{solution}{pb:b2:eruptive-styles:1}
\textbf{1.} $\ln T = 5.298$, $4.605$, $3.871$, $2.639$, $1.386$. De hellingen tussen
opeenvolgende isopachen zijn $-0.0693$, $-0.0367$, $-0.0308$ en $-0.0313$ per km: één
steil segment tussen 10 en \qty{20}{km}, een flauw van 40 tot
\qty{120}{km}; de knik ligt tussen 20 en \qty{40}{km}.

\textbf{2.} $k_1 = \ln(200/100)/10 = \qty{0.0693}{km^{-1}}$; $T_0 = 200\,e^{0.693} =
\qty{400}{cm}$.

\textbf{3.} $167e^{-1.244} = \qty{48.1}{cm}$, $167e^{-2.488} = \qty{13.9}{cm}$,
$167e^{-3.732} = \qty{4.00}{cm}$: binnen \qty{1}{\percent} van 48, 14 en 4.0.

\textbf{4.} $400e^{-0.0693X} = 167e^{-0.0311X}$ geeft $X = \ln(400/167)/(0.0693 -
0.0311) = 0.873/0.0382 = \qty{22.9}{km}$.

\textbf{5.} $\ln2/0.0693 = \qty{10.0}{km}$ en $\ln2/0.0311 = \qty{22.3}{km}$.

\textbf{6.} $T_0 = \qty{0.0040}{km}$: $V = 2\times0.0040/0.0693^2 =
\qty{1.67}{km^3}$.

\textbf{7.} $e^{-0.0693\times22.9} = 0.205$; $(0.0311\times22.9 + 1)/0.0311^2 = 1770$;
$(0.0693\times22.9 + 1)/0.0693^2 = 539$; correctie $2\times0.0040\times0.205\times1231 =
\qty{2.02}{km^3}$; $V = 1.67 + 2.02 = \qty{3.68}{km^3}$, waarvan de correctie
\qty{55}{\percent} is.

\textbf{8.} $3.68/1.67 = 2.2$. De fijne distale as is dun, gemakkelijk weggeërodeerd,
herwerkt of in zee gevallen; in een oude afzetting is ze vaak verloren of nooit in kaart
gebracht, en het volume wordt dan met een factor twee of meer onderschat.

\textbf{9.} $\qty{3.68e9}{m^3}$ ligt tussen $10^9$ en $10^{10}$: VEI~5.

\textbf{10.} $V_{\mathrm{DRE}} = 3.68\times900/2400 = \qty{1.38}{km^3}$.

\textbf{11.} $m = \num{3.68e9}\times900 = \qty{3.31e12}{kg}$.

\textbf{12.} $M = \log_{10}(\num{3.31e12}) - 7 = 12.52 - 7 = 5.52$.

\textbf{13.} $3.68/1.3 = 2.8$ keer.

\textbf{14.} Een kolom hoger dan \qty{25}{km}, en een Plinische stijl.

\textbf{15.} $Q = \num{3.31e12}/36000 = \qty{9.2e7}{kg/s}$; $I = 7.96 + 3 = 11.0$.

\textbf{16.} $0.304\times(\num{9.2e7})^{0.241} = 0.304\times83.0 = \qty{25.2}{km}$.

\textbf{17.} $0.345\times(\num{9.2e7})^{0.226} = 0.345\times63.1 = \qty{21.8}{km}$.

\textbf{18.} De gemeten \qty{31}{km} ligt \qty{23}{\percent} boven de eerste
voorspelling en \qty{42}{\percent} boven de tweede, binnen de
spreiding van \qtyrange{53}{60}{\percent} van de wetten: verenigbaar.

\textbf{19.} $Q = (31/0.304)^{1/0.241} = 102.0^{4.149} = \qty{2.2e8}{kg/s}$, 2.3 keer
het gemiddelde debiet. De klasten leggen de hoogste kolom vast, op het hoogtepunt van de
eruptie, terwijl vraag~15 middelt over de hele fase; en een deel van de massa
kan in de afzetting ontbreken (distale as, \omterm{def:b2:eruptive-styles:pdc}{pyroclastische stromen}).

\textbf{20.} Ja: 22 tot \qty{31}{km}, ruim boven de tropopauze op \qty{16}{km}.
Gas dat in de stratosfeer terechtkomt, verspreidt zich rond de aardbol; het sleutelgas is
zwaveldioxide, en de twintig miljoen ton die de Pinatubo er in 1991 bracht,
koelde de aardbol twee jaar lang met ongeveer een halve graad af.

\textbf{21.} $\lambda_3 = 399/75 = 5.32$, $\lambda_4 = 95/125 = 0.760$, $\lambda_5 = 22/225 =
0.0978$, $\lambda_6 = 8/525 = 0.0152$ per jaar; logaritmen $0.726$, $-0.119$, $-1.010$,
$-1.817$.

\textbf{22.} De rechte geeft $10^{0.73} = 5.37$, $10^{-0.12} = 0.759$, $10^{-0.97} =
0.107$ en $10^{-1.82} = 0.0151$: binnen \qty{1}{\percent}, \qty{0}{\percent},
\qty{10}{\percent} en \qty{1}{\percent} van de tellingen.

\textbf{23.} Het archief van grote erupties blijft verder terug in de tijd volledig,
omdat ze van verder worden opgemerkt en hun afzettingen langer overleven; een
venster dat in 1500 begint, zou erupties van VEI~3 sterk ondertellen, terwijl een venster
dat in 1950 begint, bijna geen eruptie van VEI~6 zou bevatten.

\textbf{24.} $\lambda_5 = \num{0.107}$ per jaar, een herhalingstijd van $1/0.107 =
\num{9.3}\text{ jaar}$; $1 - e^{-1.07} = 1 - 0.343 = \qty{66}{\percent}$ in tien jaar.

\textbf{25.} $\lambda_7 = 10^{3.28 - 5.95} = 10^{-2.67} = 0.00214$ per jaar, één om de
470 jaar; het archief geeft $2/1025 = 0.00195$ per jaar: overeenstemming binnen
\qty{10}{\percent}, voor een punt dat niet voor de fit is gebruikt.

\textbf{26.} $\lambda_6 = 10^{3.28 - 5.10} = 10^{-1.82} = 0.0151$ per jaar; $\lambda t = 1.51$
en $1 - e^{-1.51} = 1 - 0.220 = 0.780$: de kans op minstens één eruptie
van VEI~6 of meer ergens op aarde in de komende honderd jaar is
\textbf{\qty[text-series-to-math=true]{78}{\percent}}. Ze berust op een constante frequentie,
op erupties die onafhankelijk van elkaar zijn (een poissonproces), en op een
catalogus die binnen zijn vensters volledig is.
\end{solution}
